% ============================================================
%  TAYLOR & FRANCIS — Applied Operations and Analytics
%  Full journal template — single self-contained main.tex
%  Replicates: typography, column layout, header/footer,
%  coloured section headings, table style, figure style,
%  algorithm box, pairwise-matrix table, AHP weight table,
%  boxplot figure stub, citation & reference style (APA-like),
%  abstract box, keywords, article-history block.
% ============================================================

% ── Engine ──────────────────────────────────────────────────
\documentclass[10pt,twocolumn]{article}

% ── Core packages ───────────────────────────────────────────
\usepackage[T1]{fontenc}
\usepackage[utf8]{inputenc}

% ── Fonts ────────────────────────────────────────────────────
% Body: Times New Roman equivalent
\usepackage{mathptmx}          % serif body + math = Times
\usepackage[scaled=0.92]{helvet} % sans-serif (section labels) = Helvetica/Arial
\usepackage{courier}           % monospace
% Title-block sans (mimics the bold sans used in the journal header)
\renewcommand{\sfdefault}{phv}

% ── Page geometry ────────────────────────────────────────────
% Measured from the PDF: ~170 mm text width, 242 mm text height,
% two 82 mm columns with 6 mm gutter
\usepackage[
  paperwidth  = 210mm,
  paperheight = 297mm,
  top         = 22mm,
  bottom      = 28mm,
  left        = 20mm,
  right       = 20mm,
  columnsep   = 6mm,
  headsep     = 8pt,
  footskip    = 16pt
]{geometry}

% ── Colours ─────────────────────────────────────────────────
\usepackage[dvipsnames,table]{xcolor}
% Journal blue used in section headings, article title, table rules
\definecolor{tfjournalblue}{RGB}{0,80,158}
% Light header-rule grey
\definecolor{tfgrey}{RGB}{80,80,80}
% Table shading (alternating rows — not used here but defined)
\definecolor{tfshade}{RGB}{230,237,245}

% ── Header & footer ─────────────────────────────────────────
\usepackage{fancyhdr}
\pagestyle{fancy}
\fancyhf{}
% ── Mimics the journal running head exactly ─────────────────
\fancyhead[L]{\footnotesize\sffamily\color{tfgrey}APPLIED OPERATIONS AND ANALYTICS}
\fancyhead[R]{\footnotesize\sffamily\color{tfgrey}\thepage}
\fancyfoot[C]{}
\renewcommand{\headrulewidth}{0.6pt}
\renewcommand{\headrule}{\color{tfgrey}\hrule width\headwidth height\headrulewidth}
% First page uses a different style (title page style)
\fancypagestyle{firstpage}{
  \fancyhf{}
  \fancyhead[L]{\footnotesize\sffamily\color{tfgrey}RESEARCH ARTICLE}
  \fancyhead[R]{%
    \footnotesize\sffamily\color{tfjournalblue}
    \textbf{a} OPEN ACCESS \quad \textcolor{tfgrey}{$\checkmark$}%
  }
  \fancyfoot[C]{}
  \renewcommand{\headrulewidth}{0.6pt}
  \renewcommand{\headrule}{\color{tfgrey}\hrule width\headwidth height\headrulewidth}
}

% ── Spacing ──────────────────────────────────────────────────
\usepackage{setspace}
\setstretch{1.05}               % body leading — slightly loose like the journal
\setlength{\parindent}{0pt}
\setlength{\parskip}{4pt plus 1pt}

% ── Section heading style ────────────────────────────────────
\usepackage{titlesec}
% Level-1: bold italic, blue, 10.5 pt, flush left, ruled above
\titleformat{\section}
  {\normalfont\fontsize{10.5}{13}\bfseries\itshape\color{tfjournalblue}}
  {\thesection.}
  {0.5em}
  {}
  [\vspace{-2pt}\color{tfjournalblue}\rule{\columnwidth}{0.4pt}]
\titlespacing{\section}{0pt}{10pt plus 2pt}{4pt}

% Level-2: bold italic, blue, 10 pt
\titleformat{\subsection}
  {\normalfont\fontsize{10}{12}\bfseries\itshape\color{tfjournalblue}}
  {\thesubsection.}
  {0.5em}
  {}
\titlespacing{\subsection}{0pt}{8pt plus 2pt}{3pt}

% Level-3: italic only, black, run-in
\titleformat{\subsubsection}[runin]
  {\normalfont\fontsize{10}{12}\itshape}
  {\thesubsubsection.}
  {0.5em}
  {}
  [.\quad]
\titlespacing{\subsubsection}{0pt}{6pt}{0pt}

% ── Math & algorithms ────────────────────────────────────────
\usepackage{amsmath,amssymb,amsthm}
\usepackage{algorithm}
\usepackage{algorithmic}
\renewcommand{\algorithmicrequire}{\textbf{Input:}}
\renewcommand{\algorithmicensure}{\textbf{Output:}}

% ── Figures ──────────────────────────────────────────────────
\usepackage{graphicx}
\usepackage{float}
\usepackage{tikz}
\usepackage{pgfplots}
\pgfplotsset{compat=1.18}
\usetikzlibrary{arrows.meta,positioning,shapes.geometric,patterns,calc}

% ── Tables ───────────────────────────────────────────────────
\usepackage{booktabs}
\usepackage{array}
\usepackage{tabularx}
\usepackage{multirow}
% Custom table font size (journal uses ~8.5–9 pt in tables)
\newcommand{\tabfont}{\fontsize{8.5}{10.5}\selectfont}

% ── Caption style ────────────────────────────────────────────
\usepackage[
  format       = plain,
  labelfont    = {bf, small},
  textfont     = {it, small},
  labelsep     = period,
  justification= raggedright,
  singlelinecheck = false
]{caption}

% ── Lists ────────────────────────────────────────────────────
\usepackage{enumitem}
\setlist[itemize]{leftmargin=*,topsep=3pt,itemsep=1pt,parsep=0pt}
\setlist[enumerate]{leftmargin=*,topsep=3pt,itemsep=1pt,parsep=0pt}

% ── Hyperlinks & references ──────────────────────────────────
\usepackage[
  colorlinks = true,
  linkcolor  = tfjournalblue,
  citecolor  = tfjournalblue,
  urlcolor   = tfjournalblue
]{hyperref}
\usepackage{url}
% APA-style author-year citations via natbib
\usepackage[round,authoryear,sort&compress]{natbib}
\setlength{\bibsep}{3pt plus 1pt}

% ── Miscellaneous ────────────────────────────────────────────
\usepackage{microtype}   % kerning & protrusion
\usepackage{balance}     % balance last page columns
\usepackage{orcidlink}   % \orcidlink{0000-…}

% ── Utility: thin rule for abstract box ─────────────────────
\newcommand{\thinrule}{\vspace{2pt}\noindent\color{black}\rule{\columnwidth}{0.4pt}\vspace{2pt}}

% ============================================================
%  DOCUMENT
% ============================================================
\begin{document}

% ── Journal-header bar (one-column span across top) ─────────
\twocolumn[{%
%
% -------- top journal strip --------------------------------
\noindent
\begin{minipage}[t]{0.18\textwidth}
  \includegraphics[width=\linewidth]{example-image}% replace with journal cover thumbnail
\end{minipage}
\hfill
\begin{minipage}[t]{0.55\textwidth}
  \vspace{4pt}
  {\fontsize{14}{17}\selectfont\sffamily\bfseries\color{tfjournalblue}
   Applied Operations and Analytics}\\[3pt]
  {\footnotesize\sffamily\color{tfgrey}
   ISSN: 2996-6892 (Online)\quad
   Journal homepage: \url{www.tandfonline.com/journals/taoa20}}
\end{minipage}
\hfill
\begin{minipage}[t]{0.15\textwidth}
  \raggedleft
  \includegraphics[width=\linewidth]{example-image}% T&F logo
\end{minipage}

\vspace{10pt}
\noindent\color{tfgrey}\rule{\textwidth}{0.6pt}
\vspace{10pt}

% -------- article title ------------------------------------
{\fontsize{16}{20}\selectfont\sffamily\bfseries\color{tfjournalblue}
A [Self-Adjusting / Adaptive] Decision Support System for \\
Strategic [Domain] Optimization Using Integrated [Method A] \\
and [Method B]\par}

\vspace{8pt}

% -------- authors ------------------------------------------
{\fontsize{11}{14}\selectfont\sffamily\bfseries
First Author\textsuperscript{a}
\orcidlink{0000-0000-0000-0001}
\quad and \quad
Second Author\textsuperscript{a}
\orcidlink{0000-0000-0000-0002}\par}

\vspace{4pt}

{\footnotesize\itshape
\textsuperscript{a}Department of Industrial Engineering and Management Systems,
University Name, City, Country\par}

\vspace{10pt}

% -------- cite / link block --------------------------------
{\footnotesize\sffamily
\textbf{To cite this article:} First Author \& Second Author (202X)
A [Self-Adjusting] Decision Support System for [topic],
\textit{Applied Operations and Analytics}, X:X, 1--19,
DOI: 10.1080/XXXXXXXX.202X.XXXXXXX

\vspace{3pt}
\textbf{To link to this article:} \url{https://doi.org/10.1080/XXXXXXXX.202X.XXXXXXX}
\par}

\vspace{6pt}
\noindent\color{tfgrey}\rule{\textwidth}{0.4pt}
\vspace{10pt}

% -------- two-column abstract block -------------------------
\begin{minipage}[t]{0.55\textwidth}
  \setlength{\fboxsep}{6pt}%
  \setlength{\fboxrule}{0.5pt}%
  \noindent\framebox[\linewidth][l]{%
    \begin{minipage}{\dimexpr\linewidth-2\fboxsep-2\fboxrule}%
    \footnotesize
    \textbf{ABSTRACT}\\[4pt]
    This study proposes a novel [Framework Name] that integrates
    [Method A] and [Method B]. The system is designed to optimize
    [problem domain] under dynamic and uncertain conditions. The
    framework addresses a critical challenge in strategic
    decision-making by combining [static component] with a
    [dynamic component] that adaptively adjusts its actions in
    response to real-time environmental feedback. [Method A]-derived
    weights for Key Performance Indicators (KPIs) are incorporated
    into a strategically designed reward function that guides the
    [agent/algorithm], operating under an $\varepsilon$-greedy
    policy, to learn an optimal long-term [allocation/selection]
    strategy. This ensures that adaptive decisions remain aligned
    with organizational priorities while responding flexibly to
    operational changes. The framework was validated via [simulation
    platform], demonstrating a [X]\% improvement in cumulative
    [metric] compared to conventional non-adaptive approaches.
    Results indicate that the [system] consistently generates
    high-quality, data-driven decisions, enhancing [KPI 1], [KPI 2],
    and long-term [outcome]. This framework provides actionable
    insights for managers and decision-makers seeking to integrate
    strategic goal-setting with adaptive operational decision-making
    in complex and uncertain environments.
    \end{minipage}%
  }%
\end{minipage}
\hfill
\begin{minipage}[t]{0.40\textwidth}
  \footnotesize
  \textbf{ARTICLE HISTORY}\\
  Received DD Month YYYY\\
  Accepted DD Month YYYY\\[6pt]
  \textbf{KEYWORDS}\\
  [Keyword 1];
  [Keyword 2];
  [Keyword 3];
  [Keyword 4];
  [Keyword 5];
  [Keyword 6];
  [Keyword 7]
\end{minipage}

\vspace{14pt}
}]% end \twocolumn[...]

\thispagestyle{firstpage}

% ============================================================
\section{Introduction}
% ============================================================

[Topic area] is a core organizational capability, essential for
aligning investments with long-term strategic goals and achieving
organizational excellence \citep{Author2021, Author2016}. In
today's dynamic and uncertain business environment, organizational
success increasingly depends on effective and timely strategic
decision-making, often supported by robust Decision Support Systems
(DSS) \citep{Author2024, Author2021b}. This requires continuous
development of DSS and [domain] functionalities to maintain
organizational agility and respond effectively to environmental
changes \citep{Author2023, Author2023b, Author2020}.

While structured [domain] methodologies play a vital role in
establishing evaluation criteria and processes \citep{Author2020b,
Author2017}, their effectiveness is fundamentally challenged by
environmental volatility \citep{Author2016b, Author2012}. Relying
on static multi-criteria decision-making (MCDM) models (e.g.,
simple [Method A] or [Variant]) leads to strategic rigidity in
resource allocation, as the strategic weights of criteria fail to
update quickly enough with shifting operational conditions and
emerging risks \citep{Author2017b}. This situation necessitates a
critical evaluation of decision support architectures within
[context] firms. Consequently, a key management challenge arises:
How can an automated and rational Adaptive Decision Support System
(ADSS) be designed to dynamically adjust [domain] criteria in
real-time, ensuring continuous alignment with evolving
environmental conditions and strategic priorities?

Significant advancements in intelligent methods, particularly
Reinforcement Learning (RL), offer a powerful new paradigm for
sequential and dynamic decision problems \citep{Author2021c}. RL
has demonstrated strong performance and adaptability in optimizing
dynamic resource allocation in complex systems as well as in
financial portfolio management \citep{Author2025}, thereby boosting
organizational resilience. The development of hybrid DSS
architectures that integrate structured approaches like MCDM with
optimization techniques is vital for effective project selection in
complex environments \citep{Author2025b}. Nevertheless, the [domain]
literature lacks a coherent hybrid ADSS model that successfully
integrates formal strategic preference setting (i.e., managerial
KPI weighting) with dynamic operational decision-making via RL.

This research addresses this critical gap by introducing a novel
[Method B]--[Method A] hybrid framework for Adaptive Decision
Support Systems (ADSS). The core contribution lies in
synergistically integrating the [Algorithm] (for dynamic,
autonomous action selection) with the [Method A]. [Method A]
serves as a robust tool to calculate validated KPI weights,
ensuring the RL system's reward function is directly and
authentically tied to management's strategic priorities
\citep{Author2011}. This combination not only enhances decision
quality in the face of uncertainty but also provides a near-optimal
policy for [domain] management that consistently remains aligned
with the organization's long-term strategic direction.

Results from a dynamic simulation using [Platform] demonstrate
that the [Method B]--[Method A] hybrid framework significantly
enhances strategic alignment and [metric] compared to static
methods, providing [domain] managers with a self-adjusting DSS for
effective decision-making in uncertain environments.

% ============================================================
\section{Literature Review}
% ============================================================

\subsection{Theoretical Foundations and Requirements for Adaptive
Decision Support Systems (ADSS)}

In the management literature, [Domain] Management is increasingly
viewed not merely as a mechanical tool for [selection] but as a
Governance Structure and a Strategy Enactment mechanism. Effective
[domain] management enables organizations to achieve Ambidexterity
by managing the simultaneous pursuit of exploration and
exploitation \citep{Author2010}. This strategic requirement
directly translates into the necessity for Adaptive Decision
Support System (ADSS) architectures that can computationally
support dynamic strategy.

Portfolio Governance plays a key role in establishing
transparency, allocating authority, and creating clear processes
for strategic decision-making \citep{Author1999}. From a systems
perspective, Dynamic Capabilities Theory provides the necessary
theoretical justification for strategic flexibility in [domain],
emphasizing the organization's ability to sense, seize, and
transform resources to maintain agility \citep{Author2007}. These
capabilities are vital for organizational excellence and sustained
survival in fast-paced, competitive environments characterized by
technological convergence and market uncertainty \citep{Author2016c}.
The central challenge for DSS design is developing a computational
framework that can accurately and autonomously operationalize these
dynamic capabilities.

\subsection{The Limitations of Multi-Criteria Decision-Making
(MCDM) Models in Dynamic Decision Support}

Traditional decision support for resource allocation in [domain]
is primarily guided by Multi-Criteria Decision-Making (MCDM)
models. The Analytic Hierarchy Process ([Method A]), founded by
\citet{Saaty1990}, is recognized as a powerful tool for converting
subjective qualitative managerial judgments into quantitative
relative weights for decision criteria \citep{Author1990}. The key
value of [Method A] and its family (e.g., ANP) lies in their
ability to: (1) manage consistency in managerial judgments and (2)
provide transparency in the prioritization process.

However, the main structural weakness of [Method A] and its
derivatives is their inherent inability to manage the sequential
and adaptive nature of [domain] decisions \citep{Author2012b}.
[Domain] decisions span long time horizons and necessitate frequent
re-evaluation \citep{Author2007b}. Static MCDM methods offer only
a ``snapshot'' of strategic preferences at a single point in time,
failing to internalize the dynamic changes in criteria weights and
resource constraints throughout the [domain] lifecycle. This
deficiency leads to a critical breakdown in decision support
efficacy \citep{Author2023c, Author2021d}.

\subsection{Reinforcement Learning (RL) as a Computational
Solution for Sequential Decision Systems}

To address the shortcomings of static decision support models, the
scientific literature has shifted toward algorithmic frameworks
built upon Markov Decision Processes (MDPs)
\citep{Author2021e}. Reinforcement Learning (RL), rooted in
control theory and machine learning, is the most robust and
suitable algorithmic paradigm for solving sequential decision
problems in dynamic environments \citep{Sutton1998,Author2007c}.

RL, especially the Q-Learning algorithm developed by
\citet{Watkins1992}, allows the intelligent agent to learn an
optimal policy without prior knowledge of the environment model
(model-free learning) through repeated interaction and receiving a
Reward signal \citep{Author2017c}. RL applications in related
management fields are growing rapidly and prove its efficacy in
adaptive decision-making:

\begin{itemize}
  \item \textbf{Resource Management:} RL is used for optimizing
        dynamic resource allocation in complex systems,
        outperforming conventional methods \citep{Author2024b}.
  \item \textbf{Operations \& Supply Chain:} In dynamic inventory
        and supply chain management, RL enables real-time, adaptive
        control policies \citep{Author2022}.
  \item \textbf{Financial Portfolio:} RL has demonstrated strong
        performance in selecting financial portfolios under high
        market volatility and sequential risk.
\end{itemize}

Despite these advancements, a central gap persists in [domain]
research and the design of hybrid DSS architectures: existing RL
implementations in management simplify the problem by assuming
pre-defined, typically financial, reward functions
\citep{Author2022b}. Without a structured mechanism to tie
high-level strategic preferences ([Method A] output) to the
dynamic reward signal (RL input), the learning system may drift,
leading to policies that are operationally efficient but
strategically misaligned \citep{Author2023d}. This constitutes a
significant system design challenge for adaptive DSS.

\subsection{Synthesis and Research Contribution}

This research directly addresses this theoretical and practical gap
by formally integrating the prescriptive power of [Method A]
(representing static strategic thought) and the adaptive capability
of RL (representing dynamic operational decision-making)
\citep{Author2023e}.

The main contribution of this study is the creation of a novel
[Method B]--[Method A] hybrid Adaptive Decision Support System
(ADSS) framework where the [Method A] output (strategic criteria
weights) serves as a Strategic Filter in designing the Q-Learning
reward function. This dual mechanism provides the first system
that: (1) possesses the agility required to respond to
environmental changes (enhancing dynamic capabilities) and (2)
ensures that all dynamic [domain] actions are consistently aligned
with the approved strategic vision and goals of senior management
\citep{Author2021f}.

% ============================================================
\section{Methodology: [Framework Name] (ADSS) Framework}
% ============================================================

This section presents the methodological design of the proposed
[Method B] -- [Method A] Adaptive Decision Support System (ADSS).

\subsection{Overview and Formalization of the Hybrid ADSS
Architecture}

The proposed model is a [Method B]--[Method A] hybrid Adaptive
Decision Support System (ADSS) framework designed to optimize
[domain] selection under uncertainty dynamically. The architecture
formally integrates the prescriptive power of Multi-Criteria
Decision-Making (MCDM) with the adaptive capability of
Reinforcement Learning (RL) to address the limitations of static
decision-making in turbulent environments. To directly resolve
this limitation, the proposed [Method B]--[Method A] ADSS
incorporates a Strategic Filter that evaluates and constrains each
candidate RL action against the strategic hierarchy defined
through [Method A]. The sequential decision problem is modeled as
a Markov Decision Process (MDP), defined by the tuple
$\langle S, A, T, R \rangle$.

\subsection{Markov Decision Process (MDP) Formulation}

Q-learning was selected as the learning algorithm due to the fully
discrete and relatively low-dimensional state--action space of the
[domain] environment, which makes tabular methods both feasible and
efficient. Unlike DQN-based approaches, which are designed for
high-dimensional or continuous state spaces requiring function
approximation, the present setting does not necessitate neural
network estimation. Q-learning also offers an advantage over SARSA
in this application because its off-policy nature enables stable
learning of a target greedy policy while allowing stochastic
exploratory behavior during training. Furthermore, the tabular
Q-learning framework provides full transparency and
interpretability of learned values, which is essential for
managerial decision-support contexts.

\subsubsection{Q-Learning Algorithm}
In this study, we employ Reinforcement Learning (RL) to enable
adaptive decision-making in [domain] management under uncertainty.
RL is well-suited for dynamic environments, allowing agents to
learn through trial and error by interacting with the environment
and adjusting their policy accordingly \citep{Sutton1998}.

Specifically, we utilize the Q-learning algorithm, a model-free
method where the agent observes the current state (defined by
performance indicators), selects an action, receives a reward, and
iteratively updates its Q-values. The Q-Learning algorithm
determines the optimal state-action mapping through gradual
learning. Continuous adaptation is achieved via the Q-value update
rule:

\begin{equation}
  Q(s_t, a_t) \leftarrow Q(s_t, a_t) +
  \alpha \bigl[ r_t + \gamma \max_{a'} Q(S_{t+1}, a') - Q(s_t, a_t) \bigr]
  \label{eq:qupdate}
\end{equation}

\noindent where:
\begin{itemize}
  \item $\alpha$ is the learning rate $(0 < \alpha < 1)$,
        determining how quickly the algorithm adapts to new data,
  \item $\gamma$ is the discount factor $(0 < \gamma < 1)$ for
        future rewards,
  \item $R_t$ is the reward received at time step $t$,
  \item $Q(s_t, a_t)$ is the current value of the state-action
        pair $(s_t, a_t)$, and $\max Q(S_{t+1}, a')$ represents the
        maximum Q-value over all possible actions $a'$ in the next
        state $S_{t+1}$.
\end{itemize}

% ---- Algorithm box ----------------------------------------
\begin{algorithm}[H]
\caption{Q-Learning Approach}
\label{alg:qlearning}
\begin{algorithmic}[1]
  \REQUIRE State space $S$, action space $A$, learning rate
           $\alpha$, discount factor $\gamma$, exploration rate
           $\varepsilon$
  \ENSURE  Optimized Q-value function $Q(s,a)$
  \STATE Initialize Q-values: $Q(s,a)\leftarrow 0$,
         $\forall\, s\in S,\; a\in A$
  \REPEAT
    \STATE Observe current state $s = s_t$
    \STATE Select action $a_t$ using $\varepsilon$-greedy policy:
    \STATE \quad With probability $\varepsilon$, choose a random
           action (exploration)
    \STATE \quad With probability $1-\varepsilon$, choose
           $a_t = \arg\max_a Q(s,a)$ (exploitation)
    \STATE Execute action $a_t$, observe reward $r_t$, and next
           state $s' = s_{t+1}$
    \STATE Update Q-value using Equation~\eqref{eq:qupdate}
    \STATE Set $s \leftarrow s'$
  \UNTIL{convergence criterion is met
         (e.g., change in Q-values $<$ threshold)}
\end{algorithmic}
\end{algorithm}

The State Space is defined by three discrete and mutually
exclusive states: $S = \{\text{Normal},\, \text{Crisis},\,
\text{Sudden Change}\}$. These states collectively capture the
overall [domain] health and the underlying environmental
uncertainty context. The Action Space includes $N$ discrete,
[domain]-level strategic policies representing management's
responses to environmental states.

\subsection{Hybrid Reward Function Design and Strategic
Aggregation}

Unlike conventional RL frameworks with predefined reward
functions, this study introduces a hybrid reward formulation
directly grounded in managerial strategic weights derived from
[Method A]. The framework's core contribution is the Hybrid
Reward Function ($R_{\mathrm{total}}$), which directly aligns the RL
agent's learning with management's strategic preferences ([Method
A] weights). Its calculation involves two stages: Immediate KPI
Reward Estimation and Strategic Aggregation.

\subsubsection{Immediate KPI Reward Estimation}
The observed reward for a KPI, $r_{n,i}$, is based on its
percentage deviation from the strategic target. To ensure the
reward reflects stable strategic changes rather than short-term
fluctuations, a Smoothed Strategic Performance Score ($R_n$) is
calculated using a rolling window $k$:

\begin{equation}
  R_n = \frac{1}{k} \sum_{i=1}^{k} r_{n,i}
  \label{eq:smoothed}
\end{equation}

The corresponding reward structure based on action outcome is
summarized in Table~\ref{tab:rewards}.

\subsubsection{Strategic Aggregation and Temporal Weighting}
The final reward received by the agent at each time step is the
Total Portfolio Reward ($R_{\mathrm{total}}$). It is derived from
the strategically weighted aggregation of the stabilized
performance scores ($R_n$) using the [Method A]-derived strategic
weights ($w_n$):

\begin{equation}
  R_{\mathrm{total}} = \sum_{n=1}^{K} w_n R_n
  = w_1 R_1 + w_2 R_2 + \cdots + w_K R_K
  \label{eq:rtotal}
\end{equation}

\subsection{Implementation Steps of the Hybrid Framework}

The proposed decision-making framework integrates the static,
prescriptive power of [Method A] with the adaptive capability of
the Q-Learning algorithm. The implementation is executed in two
distinct phases:

\textit{Phase 1: Strategic Pre-Processing ([Method A]).}
(1) Identification and Weighting of KPIs: Define the set of the
$n$-th Key Performance Indicator ($\text{KPI}_n$) that represents
the organizational strategic goals. The relative importance
(strategic weight, $w_n$) is formally computed through [Method A],
ensuring managerial consensus and consistency. (2) State and
Action Definition: Analyze [domain] performance thresholds to
define the environmental states and the corresponding [domain]
actions available to the agent.

\textit{Phase 2: Iterative Learning (Q-Learning).}
(1) Initialization: Initialize the Q-Table $Q(s,a)\leftarrow 0$
and set the learning parameters ($\alpha, \gamma, \varepsilon$).
(2) Reward Computation: During each iteration, compute the
estimated expected reward ($R_n$) for each KPI using
Equation~\eqref{eq:smoothed}. (3) Hybrid Reward Aggregation:
Aggregate into $R_{\mathrm{total}}$ via Equation~\eqref{eq:rtotal}.
(4) Optimal Action Selection: Select $a_t$ based on $s_t$ using
the $\varepsilon$-greedy policy. (5) Q-Value Update: Execute,
observe $S_{t+1}$, and update using
Equation~\eqref{eq:qupdate}. (6) Policy Convergence: Repeat until
convergence.

\subsection{Case Study: Scenario-Based [Domain] Evaluation}

To validate the hybrid model, a scenario-based simulation was
conducted at a major [type] organization. Three representative
environmental shock scenarios were simulated. Shock~1 represents a
[Priority] Shock, modeled by increasing the [KPI] weight by 40\%
while proportionally reducing other strategic weights. Shock~2
represents an Operational Crisis Shock, modeled by reducing [KPI
A] and [KPI B] reward contributions by approximately 30\%.
Shock~3 represents a [Regulatory/Sustainability] Shock, modeled by
increasing the corresponding KPI weights by approximately 35\%.

Nine performance indicators (KPIs) were formulated:

\begin{enumerate}
  \item \textbf{KPI 1:} [Definition]
  \item \textbf{KPI 2:} [Definition]
  \item \textbf{KPI 3:} [Definition]
  \item \textbf{KPI 4:} [Definition]
  \item \textbf{KPI 5:} [Definition]
  \item \textbf{KPI 6:} [Definition]
  \item \textbf{KPI 7:} [Definition]
  \item \textbf{KPI 8:} [Definition]
  \item \textbf{KPI 9:} [Definition]
\end{enumerate}

% ============================================================
\section{Results and Discussion}
% ============================================================

\subsection{Analytical Hierarchy Process ([Method A]) for KPI
Weighting}

To quantify the relative strategic importance of each KPI, a
structured pairwise comparison method was applied using Saaty's
9-point scale. The evaluations were obtained from five senior
[domain] managers and strategy experts. Following standard [Method
A] computation, the matrix was normalized, and the mean of each
row was calculated to derive the priority vector representing the
relative weights of the nine KPIs. The resulting weights
(Table~\ref{tab:ahpweights}) indicate that [KPI 1] (0.XXXX) and
[KPI 2] (0.XXXX) dominate the strategic preference structure.

To verify the internal consistency of expert judgments, the
Consistency Index (CI) was computed:

\begin{equation}
  CI = \frac{\lambda_{\max} - n}{n - 1}
  \label{eq:ci}
\end{equation}

The Consistency Ratio (CR) was then obtained by dividing CI by the
Random Index (RI). According to Saaty's scale, the RI for $n = 9$
is 1.45:

\begin{equation}
  CR = \frac{CI}{RI}
  \label{eq:cr}
\end{equation}

A matrix is considered consistent if $CR < 0.1$. In this study,
the calculated CR satisfied this condition, confirming that the
expert evaluations were logically consistent.

\subsection{State--Action Pair Matrices for Key Indicators}

The state--action matrices for key performance indicators align
with the organization's long-term goals. Proposed actions enhance
resilience to shocks, strengthen [domain] systems, and improve
service quality, balancing short- and long-term strategy.

\subsection{Q-Learning Algorithm Implementation and Policy
Convergence}

\subsubsection{Parameter Configuration and Update Rule}
To ensure convergence stability and emphasize long-term strategic
outcomes, the learning rate ($\alpha$) was set to 0.1, and the
discount factor ($\gamma$) was set to 0.9. The $\varepsilon$-greedy
policy was implemented with an initial value of 1.0, linearly
decaying to 0.01 over the training episodes to ensure a robust
balance between exploration and exploitation.

\subsubsection{Policy Convergence and Performance Interpretation}
Convergence behavior was evaluated through the stabilization of
Q-values and the diminishing variance of cumulative rewards over
successive training episodes. The convergence patterns demonstrated
smooth and monotonic learning trajectories across all KPIs,
indicating robust policy stabilization and the absence of
oscillatory behavior. Stabilization of the total reward after
approximately 400 episodes confirms that the agent successfully
learned an optimal, steady-state strategic policy.

\subsection{Strategic Reward Aggregation and Final Reward
Calculation}

To calculate the total weighted reward, we apply
Equation~\eqref{eq:rtotal}, in which each indicator's average
reward is multiplied by its corresponding relative weight:

\begin{equation}
  R_{\mathrm{total}} = (0.3521 \times 1.556) + (0.1951 \times 2.0)
  + \cdots + (0.0272 \times 1.889) = 1.735
\end{equation}

\subsection{Ex-Post Analysis of KPI Performance Contribution}

Although the strategic priorities ($W_n$) are fixed and
exogenously determined by [Method A], it is valuable to conduct an
ex-post analysis to evaluate the relative contribution of each
KPI's performance ($R_n$) to the overall reward structure. The
relative operational contribution ($C_n$) of each indicator is
calculated as:

\begin{equation}
  C_n = \frac{R_n}{\sum_{i=1}^{9} R_i}
  \label{eq:contribution}
\end{equation}

The analysis shows [KPI X] contributed most to the total average
reward (Table~\ref{tab:contribution}), followed by [KPI Y],
indicating that the RL agent effectively managed these areas.

\subsection{Policy Convergence and Sensitivity Analysis}

This section presents simulation results from the proposed
reinforcement learning model. Figure~\ref{fig:cumreward} illustrates
the exploitation-only policy ($\varepsilon = 0$) maintaining higher
cumulative reward values compared to the exploratory policy
($\varepsilon = 0.3$) throughout the simulation. At the final
episode (Episode 200), the exploitation strategy reaches a
cumulative reward of approximately 950, whereas the exploratory
policy reaches approximately 780. This $\approx$21\% difference
reflects the expected ``exploration cost.''

Figure~\ref{fig:boxplot} presents the distribution of
$R_{\mathrm{total}}$ over 50 independent runs. The boxplot confirms
that the Hybrid [Method B]--[Method A] policy has the highest
median, the tightest spread, and the fewest extreme deviations,
supporting its robustness and repeatability.

\subsection{Comparative Performance Analysis}

The proposed hybrid [Method B]--[Method A] policy was benchmarked
against three non-adaptive control scenarios:

\begin{enumerate}
  \item \textbf{[Method A]-Only Policy (Control 1):} A static
        decision-making strategy based solely on [Method A]-derived
        priorities, without the ability to adapt to changing states.
  \item \textbf{Random Policy (Control 2):} A non-intelligent
        baseline where actions are selected randomly in every state.
  \item \textbf{Simple Weighted-Sum Heuristic (Control 3):} A
        classical multi-criteria scoring method in which each
        candidate action is evaluated using a deterministic
        weighted-sum score.
\end{enumerate}

The comparative performance metrics are presented in
Table~\ref{tab:comparison}. The results clearly demonstrate:

\begin{enumerate}
  \item \textbf{Superior Performance:} The model achieves the
        highest strategic performance ($R_{\mathrm{total}} = 1.7359$),
        confirming that the adaptive policy outperforms non-adaptive
        alternatives.
  \item \textbf{Gain over Static Policy:} The hybrid model
        improves the strategic weighted reward compared to the
        static-only policy (1.5500) by selecting optimal corrective
        actions during crisis or sudden change states.
  \item \textbf{Baseline Validation:} The random policy
        ($R_{\mathrm{total}} = 0.8500$) performs poorly, confirming
        that the learning mechanism and reward design effectively
        guide the agent toward high-performance outcomes.
\end{enumerate}

\subsection{Statistical Validation and Robustness Analysis}

To strengthen the validation, two additional analyses were
conducted: multi-run stability testing and environmental shock
robustness testing. The model was executed 50 independent times,
each for 200 episodes, under the same parameter setting
($\alpha=0.1$, $\gamma=0.9$). The aggregated statistical results
across the three policies are presented in
Table~\ref{tab:multirun}. The Hybrid [Method B]--[Method A] model
achieves the highest mean reward while also exhibiting the lowest
variability.

\subsection{Limitations and Computational Considerations}

While the proposed [Method B]--[Method A] framework is
computationally tractable for the [domain] size and KPI structure
examined in this study, the methodology may face scalability
challenges in substantially larger [domain] environments. As the
number of assets, KPIs, and admissible reallocation options
expands, the dimensionality of the state--action space grows
non-linearly. Future work may address these limitations through
state aggregation, function approximation, or migrating toward deep
reinforcement learning variants (e.g., DQN) to enhance scalability
and computational efficiency.

% ============================================================
\section{Conclusion}
% ============================================================

This study developed and validated a hybrid decision-making
framework integrating [Method A] and Q-Learning to optimize
[domain] management under uncertainty. Its core contribution is the
Strategic Hybrid Reward Function, which maps [Method A]-based
strategic priorities onto the agent's learning objectives, ensuring
adaptive and aligned policies.

Empirical validation demonstrated that the Q-values converged
steadily and the cumulative weighted rewards stabilized, confirming
the model's effectiveness in learning optimal policies focused on
high-priority indicators, particularly [KPI X]. Sensitivity
analyses on key parameters ($\alpha$, $\gamma$, $\varepsilon$)
highlighted the importance of careful tuning for stable convergence
and efficient long-term learning.

Comparative performance analysis showed that the hybrid model
achieved a higher Total Weighted Reward ($R_{\mathrm{total}}$) than
static-only policies, demonstrating its ability to learn adaptive,
superior strategies that yield greater strategic returns in dynamic
environments.

Future research should focus on:
\begin{itemize}
  \item Extending the framework toward deep reinforcement learning
        approaches capable of handling larger and more complex
        [domain] environments, including richer and
        higher-resolution state representations.
  \item Applying the model to real organizational datasets once
        such access becomes available, enabling empirical validation
        under real-world [domain] management conditions.
  \item Extending the benchmarking analysis by comparing the
        proposed hybrid [Method B]--[Method A] framework with
        advanced established [domain] optimization techniques (e.g.,
        mean--variance optimization, linear or integer programming
        models).
\end{itemize}

% ============================================================
%  TABLES
% ============================================================

% ---- Table 1: Q-Table sample ---
\begin{table}[H]
\caption{Sample Q-Table.}
\label{tab:qtable}
\tabfont
\setlength{\tabcolsep}{4pt}
\begin{tabularx}{\columnwidth}{lXXXX}
\toprule
\textbf{States} & $A_1$ & $A_2$ & $\cdots$ & $A_m$ \\
\midrule
$S_1$ & $Q_{11}$ & $Q_{12}$ & $\cdots$ & $Q_{1m}$ \\
$S_2$ & $Q_{21}$ & $Q_{22}$ & $\cdots$ & $Q_{2m}$ \\
$\vdots$ & $\vdots$ & $\vdots$ & $\ddots$ & $\vdots$ \\
$S_n$ & $Q_{n1}$ & $Q_{n2}$ & $\cdots$ & $Q_{nm}$ \\
\bottomrule
\end{tabularx}
\end{table}

% ---- Table 2: Reward assignment ---
\begin{table}[H]
\caption{Reward assignment based on action outcomes.}
\label{tab:rewards}
\tabfont
\setlength{\tabcolsep}{4pt}
\begin{tabularx}{\columnwidth}{rX}
\toprule
\textbf{Reward} & \textbf{Action Outcome} \\
\midrule
$+6$   & More than 20\% positive deviation from the target \\
$+5$   & 15\% to 20\% positive deviation \\
$+4$   & 10\% to 15\% positive deviation \\
$+3$   & 5\% to 10\% positive deviation \\
$+2$   & Up to 5\% positive deviation \\
$+1$   & No deviation from the target \\
$\phantom{+}0$  & Very slight shortfall (0\% to 1\% negative deviation) \\
$-1$   & Up to 5\% negative deviation \\
$-2$   & 5\% to 10\% negative deviation \\
$-3$   & 10\% to 15\% negative deviation \\
$-4$   & 15\% to 20\% negative deviation \\
$-5$   & More than 20\% negative deviation \\
\bottomrule
\end{tabularx}
\end{table}

% ---- Table 3: Pairwise comparison (full 9×9) ---
% Spans both columns — placed in table* environment
\begin{table*}[ht]
\caption{Pairwise comparison matrix.}
\label{tab:pairwise}
\tabfont
\setlength{\tabcolsep}{5pt}
\renewcommand{\arraystretch}{1.15}
\begin{tabular}{lrrrrrrrrr}
\toprule
 & \textbf{KPI 1} & \textbf{KPI 2} & \textbf{KPI 3}
 & \textbf{KPI 4} & \textbf{KPI 5} & \textbf{KPI 6}
 & \textbf{KPI 7} & \textbf{KPI 8} & \textbf{KPI 9} \\
\midrule
KPI 1 & 1       & 3       & 5       & 7       & 9       & 4       & 6       & 5       & 8       \\
KPI 2 & $1/3$   & 1       & 3       & 5       & 7       & 2       & 4       & 3       & 6       \\
KPI 3 & $1/5$   & $1/3$   & 1       & 3       & 5       & 1       & 2       & 2       & 4       \\
KPI 4 & $1/7$   & $1/5$   & $1/3$   & 1       & 3       & $1/2$   & 1       & $1/3$   & 2       \\
KPI 5 & $1/9$   & $1/7$   & $1/5$   & $1/3$   & 1       & $1/4$   & $1/5$   & $1/6$   & 1       \\
KPI 6 & $1/4$   & $1/2$   & 1       & 2       & 4       & 1       & 2       & 2       & 3       \\
KPI 7 & $1/6$   & $1/4$   & $1/2$   & 1       & 5       & $1/2$   & 1       & 2       & 4       \\
KPI 8 & $1/5$   & $1/3$   & $1/2$   & 3       & 6       & $1/2$   & $1/2$   & 1       & 3       \\
KPI 9 & $1/8$   & $1/6$   & $1/4$   & $1/2$   & 1       & $1/3$   & $1/4$   & $1/3$   & 1       \\
\bottomrule
\end{tabular}
\end{table*}

% ---- Table 4: AHP criteria weights ---
\begin{table}[H]
\caption{Criteria weights.}
\label{tab:ahpweights}
\tabfont
\setlength{\tabcolsep}{6pt}
\begin{tabularx}{\columnwidth}{lX}
\toprule
\textbf{Weight} & \textbf{Indicator} \\
\midrule
0.3521 & KPI 1 ([Name]) \\
0.1951 & KPI 2 ([Name]) \\
0.1047 & KPI 3 ([Name]) \\
0.0989 & KPI 4 ([Name]) \\
0.0768 & KPI 5 ([Name]) \\
0.0751 & KPI 6 ([Name]) \\
0.0475 & KPI 7 ([Name]) \\
0.0272 & KPI 8 ([Name]) \\
0.0225 & KPI 9 ([Name]) \\
\bottomrule
\end{tabularx}
\end{table}

% ---- Table 5: Operational contribution ---
\begin{table}[H]
\caption{Relative operational contribution of indicators.}
\label{tab:contribution}
\tabfont
\setlength{\tabcolsep}{6pt}
\begin{tabularx}{\columnwidth}{Xrr}
\toprule
\textbf{Indicator} & $R_n$ & $C_n$ \\
\midrule
KPI 1 & 1.556 & 0.097 \\
KPI 2 & 2.000 & 0.125 \\
KPI 3 & 1.667 & 0.104 \\
KPI 4 & 1.222 & 0.076 \\
KPI 5 & 2.111 & 0.132 \\
KPI 6 & 2.333 & 0.146 \\
KPI 7 & 1.778 & 0.111 \\
KPI 8 & 1.333 & 0.083 \\
KPI 9 & 1.889 & 0.118 \\
\bottomrule
\end{tabularx}
\end{table}

% ---- Table 6: Comparative performance ---
\begin{table*}[ht]
\caption{Comparative performance metrics for [domain] management scenarios.}
\label{tab:comparison}
\tabfont
\setlength{\tabcolsep}{5pt}
\renewcommand{\arraystretch}{1.15}
\begin{tabular}{p{3.8cm}p{7.8cm}rr}
\toprule
\textbf{Model Scenario} & \textbf{Strategy Description}
  & \textbf{$R_{\mathrm{total}}$}
  & \textbf{Gain vs.\ [A]-Only} \\
\midrule
Hybrid [B]--[A]
  & Optimal, adaptive policy learned through Q-Learning
    ($\alpha=0.1$, $\gamma=0.9$)
  & 1.7359 & $+11.99\%$ \\
Weighted-Sum Heuristic (Ctrl 3)
  & Classical scoring-based static optimisation using [A] weights
  & 1.6120 & $+4.00\%$ \\
{[A]-Only Policy (Ctrl 1)}
  & Static, non-adaptive policy (Simulated Result)
  & 1.5500 & --- \\
Random Policy (Ctrl 2)
  & Actions selected randomly in each state (Simulated Baseline)
  & 0.8500 & --- \\
\bottomrule
\end{tabular}
\end{table*}

% ---- Table 7: Multi-run validation ---
\begin{table}[H]
\caption{Multi-run statistical validation of the three decision policies.}
\label{tab:multirun}
\tabfont
\setlength{\tabcolsep}{3pt}
\begin{tabular}{lrrrrp{1.8cm}}
\toprule
\textbf{Scenario} & \textbf{Mean} & \textbf{SD} & \textbf{Var} & \textbf{CV} & \textbf{95\% CI} \\
\midrule
Hybrid [B]--[A]  & 1.732 & 0.041 & 0.00168 & 0.024 & {[1.720, 1.744]} \\
{[A]-Only}       & 1.548 & 0.079 & 0.00624 & 0.051 & {[1.526, 1.570]} \\
Random           & 0.861 & 0.133 & 0.01769 & 0.155 & {[0.824, 0.898]} \\
\bottomrule
\end{tabular}
\end{table}

% ---- Table 8: Robustness under shocks ---
\begin{table}[H]
\caption{Robustness results under environmental shocks.}
\label{tab:shocks}
\tabfont
\setlength{\tabcolsep}{6pt}
\begin{tabularx}{\columnwidth}{Xrrr}
\toprule
\textbf{Scenario} & \textbf{Mean $R_{\mathrm{total}}$}
 & \textbf{Change vs.\ Baseline} & \textbf{SD} \\
\midrule
Hybrid [Method B]--[Method A] (no shock) & 1.732 & ---      & 0.041 \\
Shock 1: [Priority] Shock                & 1.688 & $-2.54\%$ & 0.045 \\
Shock 2: Operational Crisis Shock        & 1.642 & $-5.20\%$ & 0.049 \\
Shock 3: [Regulatory] Shock             & 1.703 & $-1.68\%$ & 0.043 \\
\bottomrule
\end{tabularx}
\end{table}

% ============================================================
%  FIGURES
% ============================================================

% ---- Figure 1: Cumulative reward comparison (line chart) ---
\begin{figure}[H]
\centering
\begin{tikzpicture}
\begin{axis}[
  width          = \columnwidth,
  height         = 5.0cm,
  xlabel         = {Tick (month)},
  ylabel         = {$R_{\mathrm{total}}$ (cumulative)},
  title          = {\footnotesize$R_{\mathrm{total}}$ vs ticks (cumulative) -- Q},
  title style    = {font=\footnotesize},
  xlabel style   = {font=\footnotesize},
  ylabel style   = {font=\footnotesize},
  tick label style = {font=\scriptsize},
  legend style   = {font=\scriptsize, at={(0.97,0.05)},
                    anchor=south east, draw=none},
  xmin=0, xmax=200, ymin=0, ymax=1050,
  xtick          = {0,25,50,75,100,125,150,175,200},
  grid           = none,
  axis lines     = left,
  line width     = 0.8pt,
]
% Exploitation policy (ε=0) — steeper slope, reaches ~950
\addplot[color=tfjournalblue, thick, domain=0:200, samples=50]
  { 4.75*x };
\addlegendentry{$Q\;(\varepsilon=0)$}
% Exploratory policy (ε=0.3) — shallower, reaches ~780
\addplot[color=orange!80!black, thick, dashed, domain=0:200, samples=50]
  { 3.9*x };
\addlegendentry{$Q\;(\varepsilon=0.3)$}
\end{axis}
\end{tikzpicture}
\caption{Comparison of $R_{\mathrm{total}}$ cumulative values for
$\varepsilon = 0$ (exploitation) vs.\ $\varepsilon = 0.3$
(exploration) over the simulation horizon.}
\label{fig:cumreward}
\end{figure}

% ---- Figure 2: Total reward per episode (convergence) ------
\begin{figure}[H]
\centering
\begin{tikzpicture}
\begin{axis}[
  width          = \columnwidth,
  height         = 5.0cm,
  xlabel         = {Episodes},
  ylabel         = {Total Weighted Reward},
  title          = {\footnotesize Total Reward vs Episodes -- Q},
  title style    = {font=\footnotesize},
  xlabel style   = {font=\footnotesize},
  ylabel style   = {font=\footnotesize},
  tick label style = {font=\scriptsize},
  legend style   = {font=\scriptsize, at={(0.97,0.05)},
                    anchor=south east, draw=none},
  xmin=0, xmax=1000, ymin=2.4, ymax=3.5,
  grid           = none,
  axis lines     = left,
  line width     = 0.6pt,
]
% Noisy raw reward (simulated as random oscillation around converged mean)
\addplot[color=tfjournalblue!50, thin, domain=0:1000, samples=200]
  { 3.2 + 0.25*sin(deg(0.15*x)) * exp(-0.005*x)
    + 0.08*sin(deg(0.8*x)) };
\addlegendentry{raw}
% Moving average (orange) — converges smoothly
\addplot[color=orange!80!black, thick, domain=0:1000, samples=100]
  { 2.5 + 0.72*(1 - exp(-0.007*x)) };
\addlegendentry{moving average}
\end{axis}
\end{tikzpicture}
\caption{Total reward per episode over time.}
\label{fig:convergence}
\end{figure}

% ---- Figure 3: Boxplot — distribution of R_total -----------
\begin{figure}[H]
\centering
\begin{tikzpicture}
% Manual boxplot using TikZ — three groups
\def\bw{0.30}   % box half-width
\def\ww{0.15}   % whisker half-width

% ---- Group 1: Hybrid AHP-RL (median~1.73, IQR~1.72–1.74) --
\begin{scope}[xshift=1.5cm]
  % whisker lines
  \draw[thick] (0,1.68) -- (0,1.72);
  \draw[thick] (0,1.74) -- (0,1.79);
  \draw[thick] (-\ww,1.68) -- (\ww,1.68);
  \draw[thick] (-\ww,1.79) -- (\ww,1.79);
  % box
  \fill[fill=tfjournalblue!30] (-\bw,1.72) rectangle (\bw,1.74);
  \draw[thick]                 (-\bw,1.72) rectangle (\bw,1.74);
  % median
  \draw[very thick, orange!80!black] (-\bw,1.73) -- (\bw,1.73);
\end{scope}

% ---- Group 2: AHP-Only (median~1.55, IQR~1.53–1.57) -------
\begin{scope}[xshift=3.5cm]
  \draw[thick] (0,1.48) -- (0,1.53);
  \draw[thick] (0,1.57) -- (0,1.63);
  \draw[thick] (-\ww,1.48) -- (\ww,1.48);
  \draw[thick] (-\ww,1.63) -- (\ww,1.63);
  \fill[fill=tfjournalblue!30] (-\bw,1.53) rectangle (\bw,1.57);
  \draw[thick]                 (-\bw,1.53) rectangle (\bw,1.57);
  \draw[very thick, orange!80!black] (-\bw,1.55) -- (\bw,1.55);
  % outlier
  \node[circle,fill,inner sep=1.5pt] at (0,1.42) {};
\end{scope}

% ---- Group 3: Random (median~0.87, IQR~0.82–0.91) ---------
\begin{scope}[xshift=5.5cm]
  \draw[thick] (0,0.65) -- (0,0.82);
  \draw[thick] (0,0.91) -- (0,1.08);
  \draw[thick] (-\ww,0.65) -- (\ww,0.65);
  \draw[thick] (-\ww,1.08) -- (\ww,1.08);
  \fill[fill=tfjournalblue!30] (-\bw,0.82) rectangle (\bw,0.91);
  \draw[thick]                 (-\bw,0.82) rectangle (\bw,0.91);
  \draw[very thick, orange!80!black] (-\bw,0.87) -- (\bw,0.87);
\end{scope}

% Axes
\draw[->] (0.8,0.55) -- (0.8,1.92) node[above,font=\scriptsize]
  {Total Strategic Reward ($R_{\mathrm{total}}$)};
\draw[->] (0.8,0.55) -- (6.7,0.55);
% y ticks
\foreach \y/\label in {0.6/0.6, 0.8/0.8, 1.0/1.0,
                        1.2/1.2, 1.4/1.4, 1.6/1.6, 1.8/1.8}{
  \draw (0.75,\y) -- (0.85,\y);
  \node[left,font=\scriptsize] at (0.75,\y) {\label};
}
% x labels
\node[below,font=\footnotesize] at (1.5,0.55) {Hybrid [Method B]--[Method A]};
\node[below,font=\footnotesize] at (3.5,0.55) {[Method A]-Only};
\node[below,font=\footnotesize] at (5.5,0.55) {Random Policy};
\end{tikzpicture}
\caption{Boxplot of $R_{\mathrm{total}}$ over 50 runs for the three policies.}
\label{fig:boxplot}
\end{figure}

% ============================================================
%  AUTHOR CONTRIBUTIONS
% ============================================================
\vspace{6pt}
\noindent\textbf{Author contributions}

\noindent CRediT: \textbf{First Author:} Data curation,
Methodology, Validation, Visualization, Writing -- original draft;
\textbf{Second Author:} Investigation, Methodology, Project
administration, Supervision, Validation, Writing -- review \&
editing.

% ============================================================
%  DISCLOSURE
% ============================================================
\vspace{6pt}
\noindent\textbf{Disclosure statement}

\noindent No potential conflict of interest was reported by the
author(s).

% ============================================================
%  ORCID
% ============================================================
\vspace{6pt}
\noindent\textbf{ORCID}

\noindent First Author \orcidlink{0000-0000-0000-0001}
\url{http://orcid.org/0000-0000-0000-0001}\\
Second Author \orcidlink{0000-0000-0000-0002}
\url{http://orcid.org/0000-0000-0000-0002}

% ============================================================
%  DATA AVAILABILITY
% ============================================================
\vspace{6pt}
\noindent\textbf{Data availability statement}

\noindent The data supporting the findings of this study were
generated through simulation experiments. The simulation code and
related materials are openly available in Zenodo at
\url{https://doi.org/10.5281/zenodo.XXXXXXX}.

% ============================================================
%  REFERENCES — APA / Taylor & Francis author-year style
% ============================================================
\balance
\vspace{6pt}

% Inline bibliography (no .bib file needed)
\begin{thebibliography}{99}
\setlength{\itemsep}{2pt}
\setlength{\parskip}{0pt}
\bibitem[Author(2021)]{Author2021}
Author, A.\ B. (2021).
Title of the article: A theoretical review.
\textit{The International Journal of Business \& Management},
\textit{9}(1). \url{https://doi.org/10.XXXX/XXXXX}

\bibitem[Author et al.(2016)]{Author2016}
Author, C.\ D., Grobbelaar, S.\ S., \& Bam, W. (2016).
Exploring the link between [topic] and company success:
An empirical framework.
\textit{South African Journal of Industrial Engineering},
\textit{27}(3), 238--250.
\url{https://doi.org/10.7166/27-3-1635}

\bibitem[Author \& Carvalho(2024)]{Author2024}
Author, M.\ T., \& Carvalho, M.\ M. (2024).
Toward a multilevel framework of dynamic capabilities through
the lens of project management.
\textit{Project Management Journal}, \textit{55}(3), 247--263.
\url{https://doi.org/10.1177/87569728231207028}

\bibitem[Author et al.(2023)]{Author2023}
Author, J., Kaufmann, C., \& Kock, A. (2023).
The interplay between dynamic capabilities' dimensions and
their relationship to [domain] agility and success.
\textit{International Journal of Project Management},
\textit{41}(4), 102469.
\url{https://doi.org/10.1016/j.ijproman.2023.102469}

\bibitem[Saaty(1990)]{Saaty1990}
Saaty, T.\ L. (1990).
\textit{The analytic hierarchy process: Planning, priority
setting, resource allocation}.
McGraw-Hill.

\bibitem[Sutton \& Barto(1998)]{Sutton1998}
Sutton, R.\ S., \& Barto, A.\ G. (1998).
\textit{Reinforcement learning: An introduction} (1st ed.).
MIT Press.

\bibitem[Watkins \& Dayan(1992)]{Watkins1992}
Watkins, C.\ J.\ C.\ H., \& Dayan, P. (1992).
Q-learning.
\textit{Machine Learning}, \textit{8}(3--4), 279--292.
\url{https://doi.org/10.1023/A:1022676722315}

\bibitem[Author(2022)]{Author2022}
Author, W.\ B. (2022).
\textit{Sequential decision analytics and modeling: Modeling
with Python}.
Publisher.

\bibitem[Author(2023)]{Author2023d}
Author, V. (2023).
Hybrid fuzzy [Method A]--[Method B] approach to prioritizing
solutions for inverse reinforcement learning.
\textit{Complex \& Intelligent Systems}, \textit{9}(1), 493--513.
\url{https://doi.org/10.1007/s40747-022-00807-5}

\bibitem[Author et al.(2025)]{Author2025}
Author, R., Xi, Y., Stefanidis, A., Jiang, Z., \& Su, J. (2025).
A novel multi-agent dynamic portfolio optimization learning
system based on hierarchical deep reinforcement learning.
\textit{Complex \& Intelligent Systems}, \textit{11}(7), 311.
\url{https://doi.org/10.1007/s40747-025-01884-y}
\end{thebibliography}

\end{document}
